% Chapitre 21 — Rendu et Server Components : wasted-subtree-render, server-component-hook
% Sources : notes/11 (§4.6, §4.7, §6 ex. 6, 9, 10, 16, §8), notes/08 (§3.9-3.10, §4.8-4.10, §6.3, §8.1),
% notes/01 (§3.5, §4.9), notes/13 (§5.3, §6.8). Commit photographié : e67b10a.
% Exemples rejoués sous /tmp/ch21/ avec target/debug/reactant construit depuis e67b10a.

\chapter{Rendu et Server Components : wasted-subtree-render, server-component-hook}
\label{chap:regles-rendu-rsc}

\begin{objectifs}
  \item Comprendre pourquoi deux règles du catalogue ne lisent presque rien du domaine des valeurs.
  \item Suivre \regle{wasted-subtree-render} de bout en bout : déclencheurs, fréquence, arbre d'éléments, compte des rendus, options.
  \item Savoir lire une affirmation d'\emph{absence} (\og aucune entrée ne dépend de l'état\fg{}) et la règle de soundness qui l'accompagne : tout inconnu est un usage.
  \item Suivre \regle{server-component-hook} : faits de module, graphe serveur de Next.js, frontière \code{"use client"}, preuves de hook client.
  \item Savoir expliquer pourquoi ces deux règles plafonnent à \Warning{}, et ce que chacune refuse de conclure.
\end{objectifs}

\prerequis{la dépendance de rendu et l'arbre d'éléments (\cref{chap:inter-composants}, \cref{def:render-dep}) ;
la surface typée des règles et le sceau de sévérité (\cref{chap:api-regles}, \cref{def:certified}) ;
la détection et les faits de module (\cref{chap:frontend-detection}) ; les projets Next.js vus par le driver
(\cref{chap:projet-cli-driver}). Les relations \relation{slot_writers} et \relation{registrations}
(\cref{chap:relations}) servent à la première règle.}

% ---------------------------------------------------------------------------
\section{Deux règles à côté du domaine des valeurs}
\label{sec:regles-rendu-rsc-apercu}

Les règles des chapitres précédents interrogent toutes, d'une façon ou d'une autre, le résultat du
point fixe : la valeur abstraite d'un slot, la stabilité d'une référence, les lignes
\code{SlotWriter} d'un setter. Les deux règles de ce chapitre sont différentes. Elles posent des
questions que le domaine des valeurs ne sait pas trancher, et elles s'appuient sur des faits
\emph{calculés à côté} de lui.

\begin{itemize}
  \item \regle{wasted-subtree-render} demande : \og quand cet état change, quels sous-arbres
        re-rendent alors qu'aucune de leurs entrées n'en dépend ?\fg{} La réponse vient d'une
        analyse de flot séparée, la \emph{dépendance de rendu} (\adr{41}), composée sur l'arbre
        d'éléments du programme.
  \item \regle{server-component-hook} demande : \og ce hook est-il appelé dans un module que
        Next.js rend côté serveur ?\fg{} La réponse vient de deux faits qui ne sont pas du tout
        abstraits : une convention de noms de fichiers et le graphe des imports (\adr{26}).
\end{itemize}

Le \cref{tab:regles-rendu-rsc-apercu} résume ce qui les rapproche et ce qui les oppose. Le point
commun le plus important est le niveau : aucune des deux ne peut émettre une \Error{}. Le point
de divergence le plus instructif est la \emph{forme} de l'affirmation. La première affirme une
\emph{absence} (\og aucune entrée ne dépend de l'état\fg{}) ; la seconde affirme une
\emph{présence} (\og ce module est serveur, et il appelle un hook client\fg{}). Nous verrons que
ces deux formes n'imposent pas la même discipline de soundness.

\begin{table}[htbp]\centering\small
\begin{tabularx}{\linewidth}{>{\raggedright\arraybackslash}p{3.1cm}XX}
\toprule
 & \regle{wasted-subtree-render} & \regle{server-component-hook} \\
\midrule
Bug visé & rendus répétés de sous-arbres dont la sortie ne change pas & hook appelé dans un Server Component : le rendu lève une erreur \\
Faits consommés & \code{render_deps} (par composant), \code{RenderIndex}, \relation{slot_writers}, \relation{registrations}, \code{state_store} & \code{ModuleTable} (directives, arêtes d'import), \code{HookEntry}, \code{HookProvenance} \\
Forme de l'affirmation & absence d'usage & présence d'un fait de convention \\
Règle de prudence & tout inconnu est un usage & rien n'est serveur sans \code{"use client"} dans le programme ; seuls les hooks prouvés client comptent \\
Niveau & \Warning{} (rendus certains, coût incertain) & \Warning{} (faits hors du domaine) \\
Options & \code{minWastedRenders}, \code{continuousOnly} & aucune \\
\code{safe_check} & aucun & sur les Server Components seulement \\
Fichier & \file{src/rules/impls/wasted_subtree_render.rs} (333 l.) & \file{src/rules/impls/server_component_hook.rs} (308 l.) \\
Aides & \file{src/rules/helpers/render_tree.rs}, \file{src/engine/render_deps.rs} & \file{src/project/nextjs.rs}, \file{src/ir/module.rs} \\
Tests & \file{tests/wasted_subtree_render.rs} (13) & 5 tests unitaires, \file{tests/nextjs_project.rs} \\
\bottomrule
\end{tabularx}
\caption{Les deux règles du chapitre en un coup d'œil.}
\label{tab:regles-rendu-rsc-apercu}
\end{table}

Le chapitre suit l'ordre de difficulté croissante pour la première règle (le problème, les faits,
l'algorithme, les options, les cas limites, la soundness), puis traite la seconde, plus courte
mais plus délicate à justifier : c'est l'exemple type d'une règle qui \emph{sait} quelque chose
de vrai sans pouvoir le prouver dans le domaine abstrait.

% ===========================================================================
\section{Le rendu gaspillé : le problème}
\label{sec:regles-rendu-rsc-cascade}

\begin{react}[re-rendu en cascade]
Quand un état change, React re-rend le composant qui le possède, puis \emph{tous les éléments de
composant que ce rendu construit}, récursivement, que leurs props aient changé ou non. La
cascade ne s'arrête qu'à trois endroits :
\begin{enumerate}
  \item un élément qui n'a pas été reconstruit par ce rendu : un \code{children} reçu d'un parent
        garde son identité d'objet, React le reconnaît et saute son rendu ;
  \item un composant enveloppé dans \code{memo} dont toutes les props sont
        \code{Object.is}-égales aux précédentes ;
  \item la frontière du composant propriétaire : si l'état et les éléments qui le lisent vivent
        dans un composant plus petit, le reste de l'arbre n'est pas re-rendu du tout.
\end{enumerate}
Les écritures faites par un même handler sont regroupées en un seul rendu (batching), et une
écriture d'une valeur \code{Object.is}-égale à la valeur courante est abandonnée sans rendu
(bail-out). Voir \cref{chap:react-modele} pour le modèle complet.
\end{react}

Le cas qui coûte cher est celui d'un état écrit \emph{très souvent} : à chaque frappe dans un
champ texte, à chaque mouvement de souris, à chaque pas de défilement, à chaque tick de minuterie.
Si le composant qui le possède construit aussi un sous-arbre lourd qui n'en dépend pas, ce
sous-arbre produit la même sortie des dizaines de fois par seconde.

\begin{exemple}{Un champ de recherche à côté d'une barre latérale}{rendu-recherche}
Le fichier \file{/tmp/ch21/search.tsx} :
\begin{lstlisting}[language=TypeScript,caption={Un état de saisie au-dessus d'un sous-arbre qui ne le lit pas},label={lst:rendu-recherche}]
import { useState } from "react";

function Results({ query }: { query: string }) {
  return <p>results for {query}</p>;
}
function Sidebar() {
  return <nav>{["a", "b", "c"].map((k) => <Link key={k} to={k} />)}</nav>;
}
function Link({ to }: { to: string }) { return <a href={to}>{to}</a>; }

export default function SearchPage() {
  const [query, setQuery] = useState("");
  return (
    <main>
      <input value={query} onChange={(e) => setQuery(e.target.value)} />
      <Results query={query} />
      <Sidebar />
    </main>
  );
}
\end{lstlisting}
Chaque frappe écrit \code{query}. \code{SearchPage} re-rend ; \code{<Results>} re-rend à juste
titre, sa prop \code{query} a changé ; \code{<Sidebar>} et ses trois \code{<Link>} re-rendent pour
rien. La commande
\code{cargo run -q -- check /tmp/ch21/search.tsx --rule wasted-subtree-render --trace} donne :
\begin{sortie}
  SearchPage  (2 hooks)  /tmp/ch21/search.tsx
    warn   wasted-subtree-render  [hook:0]  (line 15:27)  each `change` event writes state `query`
      and re-renders `<Sidebar>` (at least 2 component renders, including a list) although none of
      its inputs depends on it. Move `query` and the elements that use it into their own component,
      or build the unrelated elements higher up and pass them in as `children`
       → `<Sidebar>` re-renders (at least 2 component renders, including a list) although none of
         its props depends on it (line 17:6)
\end{sortie}
(Les lignes longues sont repliées ici ; l'analyseur les imprime sur une seule ligne.) Le finding
pointe le handler \code{onChange} (ligne 15, colonne 27), désigne l'élément gaspillé
(\code{<Sidebar>}, ligne 17) et donne un coût : \emph{au moins} deux rendus de composant,
\code{Sidebar} et un \code{Link} compté une fois pour toute la liste.
\end{exemple}

Trois difficultés se cachent derrière ce diagnostic d'apparence simple.

\paragraph{Le domaine des valeurs ne répond pas à la question.} La question est \og la sortie de
\code{<Sidebar>} peut-elle dépendre de \code{query} ?\fg{}. Le fait le plus proche dans le
domaine est la stabilité versionnée d'\adr{17} (\code{Versioned(S)}, \cref{def:stabilite}) : il
dit de quelles écritures d'état une \emph{référence} peut provenir. Mais il ne vit que dans la
composante \code{reference} du produit \code{StateValue} : un booléen \code{open}, un nombre
\code{text.length} n'y portent rien, et un objet alloué (\code{\{ a: text \}}) est
\code{PerRender} quelle que soit son origine (\adr{41}, contexte). La dépendance est un fait
différent de la valeur ; \adr{41} §1 en tire une analyse séparée
(\cref{sec:regles-rendu-rsc-deps}).

\paragraph{La question est inter-composants.} Savoir que \code{SearchPage} passe \code{query} à
\code{<Results>} ne suffit pas : il faut savoir que \code{Results} l'\emph{utilise}, que
\code{Sidebar} ne reçoit rien, et que rien, sous \code{Sidebar}, ne lit \code{query} par un autre
canal (un contexte, une variable de module). Il faut donc composer des résumés par composant le
long de l'arbre d'éléments.

\paragraph{L'affirmation est une absence.} Un finding dit \og aucune entrée de \code{<Sidebar>}
ne dépend de \code{query}\fg{}. Une dépendance oubliée par l'analyse produirait donc un
\emph{faux positif}, pas un faux négatif. C'est l'inverse de la plupart des règles, où une
sur-approximation fait tirer davantage. Toute la conception de la règle découle de ce
renversement de polarité (\cref{sec:regles-rendu-rsc-soundness}).

\section{Ce que la règle consomme : dépendance de rendu et arbre d'éléments}
\label{sec:regles-rendu-rsc-deps}

La règle ne calcule presque rien elle-même : elle interroge un index programme, le
\code{RenderIndex}, bâti sur les résumés de dépendance de rendu de chaque composant. Ces deux
objets sont exposés en détail au \cref{chap:inter-composants} (\cref{def:render-dep}) ; on en
rappelle ici ce que la règle lit, dans l'ordre où elle le lit.

\subsection{Le résumé d'un composant}

\file{src/engine/render_deps.rs} exécute, sur le CFG convergé d'un composant (hooks personnalisés
et utilitaires déjà greffés), une analyse avant qui associe à chaque variable un
\emph{may-ensemble} de sources. Une \code{Source} (\src{src/engine/render_deps.rs}{53}{84}) est
une entrée du rendu \emph{dans le repère du composant} : \code{Slot(l)} (la valeur d'un slot),
\code{Setter(l)} (son setter, une capacité d'écriture qui ne change jamais),
\code{Prop(name)} et \code{AllProps}, \code{Ref(l)}, \code{Hook(l)} (le résultat d'un hook non
modélisé), \code{Context(id)} (un contexte prouvé) et \code{Module(var)} (un nom non lié ici que
quelque fonction du programme écrit). Seules les deux dernières sont indépendantes du repère :
elles traversent l'arbre d'éléments sans traduction.

Les ensembles se joignent par union ; le treillis est fini, donc pas de widening, et le plafond
de 64 tours dégrade le résumé en top, jamais en plus petit (\adr{41} §1). Une \code{Deps}
(\src{src/engine/render_deps.rs}{119}{131}) porte l'ensemble \code{set}, un drapeau \code{top},
une partie \code{gated} (les sources qu'une valeur fonction n'atteint que derrière un test de ses
propres arguments, fait \emph{must}) et \code{writes} (les noms de module qu'appeler la valeur
peut écrire). La question que pose une règle se formule comme une \code{Relevance}
(\src{src/engine/render_deps.rs}{207}{214}) : un ensemble de sources dans le repère d'un
composant, plus un drapeau \code{any_prop} (\og n'importe quelle prop\fg{}, après un spread). La
réponse élémentaire est \code{Deps::touches} (\src{src/engine/render_deps.rs}{182}{205}, détaillée
au \cref{chap:inter-composants}) : une dépendance top touche toute question non vide ;
\code{AllProps} (l'objet props pris en entier) touche toute prop nommée ; une question
\code{any_prop} touche toute prop ; sinon il faut une source commune.

Le résumé d'un composant, \code{RenderDeps}, sépare deux choses que tout le reste du chapitre
distingue :

\rustsnippet{src/engine/render_deps.rs}{301}{322}{le résumé : ce qui est utilisé, ce qui est construit}

\begin{itemize}
  \item \code{genuine} rassemble ce que le composant \emph{utilise} : sa sortie hôte (y compris
        les closures de handlers et les éléments hôtes qu'il passe en \code{children}), ses effets,
        ses appels en position d'instruction, ses écritures de membres, les arguments de ses hooks
        opaques, et les conditions et \emph{types} des éléments qu'il rend.
  \item \code{sites} liste chaque élément de composant qu'il construit, avec, prop par prop, ce
        dont la valeur peut dépendre. Une prop est \emph{transmise}, pas utilisée : c'est pourquoi un
        élément de composant s'évalue à l'ensemble vide (\code{CompApp} $\mapsto \emptyset$,
        \cref{chap:inter-composants}).
  \item \code{handlers} liste les props d'événement des éléments hôtes (\code{onChange} sur un
        \code{<input>}\dots) avec la balise, l'attribut \code{type} littéral et les dépendances de
        la closure : c'est là qu'un setter \og atterrit\fg{}.
  \item \code{effect_writes} donne, par effet, les noms de module que son corps peut écrire.
\end{itemize}

Chaque élément de composant est un \code{ElementSite} :

\rustsnippet{src/engine/render_deps.rs}{252}{277}{un élément de composant construit par le rendu}

Quatre champs comptent pour la règle : \code{guard} (les conditions sous lesquelles l'élément est
construit), \code{props} et \code{spread} (ce que ses entrées peuvent porter), \code{in_list}
(construit dans un callback passé à un appel, typiquement \code{.map} : zéro, une ou plusieurs
instances) et \code{parent} (l'élément de composant dans lequel celui-ci est imbriqué comme prop
ou enfant, ce qui permet de suivre un provider de contexte jusqu'aux éléments qu'il enveloppe).

\subsection{L'index programme}

\code{RenderIndex::build} (\src{src/rules/helpers/render_tree.rs}{117}{137}) calcule d'abord
l'ensemble programme des noms écrits (un nom libre n'est un canal que si quelque chose l'écrit),
puis un résumé \code{render_deps} par composant, puis le nombre de sites qui nomment chaque
composant. L'index est construit une seule fois par exécution, à la première demande
(\code{ProgramCache::render}, \src{src/rules/api/cache.rs}{70}{76}, un \code{OnceLock}) et partagé
par \regle{wasted-subtree-render} et \regle{state-lifted-too-high} (\cref{chap:infinite-loop}).

Passer d'un site à un composant se fait par \code{resolve}, qui n'accepte qu'une origine prouvée :

\rustsnippet[label={lst:rendu-resolve}]{src/rules/helpers/render_tree.rs}{766}{777}{un élément ne se résout que par son origine prouvée}

Le moteur de valeurs, lui, accepte aussi un nom unique (\code{resolve_child}, \adr{40}). Les deux
choix sont prudents, chacun dans sa polarité : l'inliner qui retrouverait un enfant par son seul
nom analyserait un corps de plus (au pire du travail inutile) ; la règle de cascade qui se
tromperait d'enfant \emph{cacherait} les usages du vrai.

\subsection{L'arbre d'éléments de \code{typing.tsx}}

La fixture de référence de la règle est \file{tests/fixtures/wasted_subtree_render/typing.tsx},
mesurée à l'exécution par l'oracle \file{scripts/rerender-bench/} (React 19, jsdom, en-tête de
\file{tests/wasted_subtree_render.rs}) :

\tsxsnippet[label={lst:rendu-typing}]{tests/fixtures/wasted_subtree_render/typing.tsx}{5}{20}{un champ texte à côté d'un sous-arbre coûteux}

Le résumé de \code{App} (relevé par le dossier de notes 08, §6.3, sur une sonde de débogage) a
\code{genuine = \{Slot(0), Setter(0)\}} (l'\code{<input>} lit \code{text} et sa closure capture
\code{setText}), un handler \code{change} sur \code{<input>} sans \code{type} dont les dépendances
sont \code{\{Setter(0)\}}, et deux sites : \code{<Preview>} avec \code{text = \{Slot(0)\}}, et
\code{<ExpensiveTree>} sans prop ni garde. Le résumé d'\code{ExpensiveTree} a un seul site,
\code{<Row>}, marqué \code{in_list}. La \cref{fig:rendu-arbre-typing} montre comment la question
\code{Relevance \{Slot(0)\}} se propage sur cet arbre.

\begin{figure}[htbp]\centering
\begin{tikzpicture}[
    level distance=17mm, sibling distance=38mm,
    comp/.style={bloc,minimum width=26mm},
    touche/.style={comp,draw=ra-amber,fill=white,very thick},
    gaspi/.style={comp,draw=ra-blue,fill=ra-blue!12,very thick},
    hote/.style={cfgbloc,minimum width=24mm},
    edge from parent/.style={draw=ra-gray,thick}]
  \node[comp] (app) {\code{App}\\[-1pt]{\scriptsize possède \code{text} = \code{Slot(0)}}}
    child { node[hote] (in) {\code{<input onChange>}\\ deps \code{\{Setter(0)\}}} }
    child { node[touche] (pre) {\code{<Preview>}\\[-1pt]{\scriptsize\code{text} : \code{\{Slot(0)\}}}} }
    child { node[gaspi] (exp) {\code{<ExpensiveTree>}\\[-1pt]{\scriptsize props $\emptyset$, garde $\emptyset$}}
      child { node[gaspi] (row) {\code{<Row>} \code{in_list}\\[-1pt]{\scriptsize compté une fois}} } };
  \node[etiquette,below=2mm of in,align=center] {landing \code{change}\\ \code{<input>} texte : continu};
  \node[etiquette,below=2mm of pre,align=center,text=ra-amber] {touché : une prop\\ porte \code{Slot(0)}};
  \node[etiquette,right=2mm of exp,align=left,text=ra-blue] {non touché,\\ résolu : gaspillé};
  \node[etiquette,right=2mm of row,align=left] {\code{renders} = 2,\\ \code{list} = vrai};
  \draw[flechemay] (in.north west) to[bend left=35] node[etiquette,left,pos=.5] {écrit} (app.west);
\end{tikzpicture}
\caption{L'arbre d'éléments de \code{typing.tsx} et la question \code{Relevance \{Slot(0)\}}.
Le handler de l'\code{<input>} écrit le slot (flèche pointillée) ; \code{<Preview>} est touché
(bord orangé) ; \code{<ExpensiveTree>} ne l'est pas, se résout, et coûte son propre rendu plus un
\code{<Row>} compté une fois pour toute la liste (bleu).}
\label{fig:rendu-arbre-typing}
\end{figure}

La propagation s'arrête dès qu'un élément est touché : \code{<Preview>} n'est pas un candidat, et
ce qui serait imbriqué dans lui ne le serait pas davantage (\code{parent}). Elle s'arrête aussi,
dans l'autre sens, dès qu'un élément est \emph{inconnu}. C'est l'invariant qui gouverne tout le
fichier \file{render_tree.rs}, énoncé dès son en-tête (\src{src/rules/helpers/render_tree.rs}{11}{13}) :

\begin{invariant}{Tout inconnu est un usage}{tout-inconnu-usage}
Dans les règles de cascade, l'absence d'usage est une preuve. Donc tout ce que l'analyse ne voit
pas compte comme un usage : un élément dont le composant ne se résout pas (bibliothèque, enveloppe
\code{memo} ou \code{forwardRef}, import non résolu, nom ambigu), un composant ré-entré par
récursion, la profondeur au-delà de \code{MAX_DEPTH = 64}, un résumé plafonné à top, un contexte
que l'analyse ne sait pas nommer (\code{any_context}), un écrivain dont la racine écrite est top
(\adr{41} §2).
\end{invariant}

\section{L'algorithme de wasted-subtree-render pas à pas}
\label{sec:regles-rendu-rsc-algo}

La méthode \code{check} de \file{src/rules/impls/wasted_subtree_render.rs} tient en deux cents
lignes (\src{src/rules/impls/wasted_subtree_render.rs}{60}{260}). Elle s'exécute une fois par
composant, qui joue le rôle de \emph{propriétaire} (\code{owner}) des états qu'il déclare. Son
squelette est donné par l'\cref{alg:rendu-wasted} ; on en suit les étapes dans l'ordre.

\begin{algorithm}[htbp]
\caption{\code{WastedSubtreeRender::check}, schématisé}\label{alg:rendu-wasted}
\KwIn{un composant propriétaire $o$, l'index \code{RenderIndex}, les options $m$ (\code{minWastedRenders}) et $c$ (\code{continuousOnly})}
\KwOut{au plus un \Warning{} par déclencheur}
$T \leftarrow \emptyset$ \tcp*{déclencheurs, groupés par clé}
\ForEach{slot $l$ d'un \code{HookEntry::State} de $o$}{
  \ForEach{\code{Landing} $h$ de \code{index.landings}$(o, l)$}{ ajouter $l$ au déclencheur \code{Handler}$(h)$ de $T$ }
}
\ForEach{écrivain local $w$ de \relation{slot_writers} dans un effet $e$}{
  ajouter $w.\mathit{slot}$ au déclencheur \code{Effect}$(e)$, avec les enregistrements de $e$ qui peuvent appeler le setter
}
\ForEach{déclencheur $t \in T$}{
  \lIf{$t.\mathit{writes}$ est top}{passer au suivant}
  \lIf{aucun événement, ou $c$ et aucun événement continu d'un slot à valeurs nombreuses}{passer}
  $\mathit{rel} \leftarrow \{\code{Slot}(l) \mid l \in t.\mathit{slots}\} \cup \{\code{Module}(x) \mid x \in t.\mathit{writes}\}$\;
  $W \leftarrow \code{index.wasted_siblings}(o, \mathit{rel})$\;
  \lIf{$W = \emptyset$, ou $\sum_{w\in W} w.\mathit{renders} < m$ sans liste}{passer}
  émettre un \Warning{} avec une note \code{Step::Rerender} par élément de $W$\;
}
\end{algorithm}

\subsection{Étape 1 : collecter les déclencheurs}

\begin{react}[batching]
Toutes les écritures d'état faites pendant un même handler d'événement, ou pendant un même
callback de minuterie, sont regroupées par React en \emph{un seul} rendu. Deux setters appelés
l'un après l'autre dans un \code{onChange} ne coûtent qu'un rendu du propriétaire.
\end{react}

La règle ne raisonne donc pas slot par slot mais \terme{déclencheur} par déclencheur : un handler
d'événement, ou l'ensemble des listeners et minuteries qu'un effet enregistre. Un déclencheur
écrit un \emph{ensemble} de slots, qui re-rendent ensemble. La clé qui décide que deux écritures
appartiennent au même lot est \code{TriggerKey}, et ce qu'on accumule pour chaque clé est un
\code{Trigger} :

\rustsnippet[label={lst:rendu-trigger}]{src/rules/impls/wasted_subtree_render.rs}{263}{292}{ce qui fait de deux écritures un seul lot}

Un déclencheur \code{Handler} est identifié par quatre choses : l'élément du propriétaire par
lequel le setter sort (\code{via}), le composant qui construit l'élément hôte portant le handler,
le span de la prop de handler et le nom de l'événement. Deux instances d'un même enfant sont donc
deux déclencheurs, parce que leurs \code{via} diffèrent. Un déclencheur \code{Effect} est
identifié par le label de l'effet.

\paragraph{Les handlers : où le setter atterrit.} Un setter n'est pas forcément appelé dans le
propriétaire : il peut être passé en prop, traverser plusieurs composants, et n'être appelé que
par un \code{<input>} trois niveaux plus bas. \code{RenderIndex::landings} répond à la question
\og où ce setter peut-il être appelé par un événement ?\fg{} en posant la question
\code{Relevance \{Setter(l)\}} et en descendant l'arbre d'éléments. Chaque réponse est un
\code{Landing} :

\rustsnippet{src/rules/helpers/render_tree.rs}{77}{97}{l'endroit où une capacité d'écriture est appelée}

La descente (\code{land}, \src{src/rules/helpers/render_tree.rs}{406}{512}) examine d'abord les
handlers hôtes du composant courant dont les dépendances touchent la question : chacun donne un
atterrissage, avec la balise et l'attribut \code{type}. Elle suit ensuite chaque prop d'un
élément de composant qui porte la capacité, en traduisant la question dans le repère de l'enfant
(\code{Setter(0)} du parent devient \code{Prop(onChange)} de l'enfant). Si une prop
d'événement (\code{onX}) entre dans un composant sans atterrir plus bas (composant opaque, enfant
qui ne l'appelle que dans un effet), l'atterrissage est posé sur l'élément lui-même, avec une
cible \code{HandlerTarget::Component}. Au passage, la descente accumule \code{keyed} (un saut
n'appelle la capacité que derrière un test de l'argument d'événement) et \code{writes} (les noms
de module que les closures traversées peuvent écrire).

\begin{exemple}{Un setter passé à un enfant}{rendu-child-setter}
La fixture \file{tests/fixtures/wasted_subtree_render/child_setter.tsx} :
\tsxsnippet{tests/fixtures/wasted_subtree_render/child_setter.tsx}{5}{22}{le setter descend dans \code{<Field>}}
\begin{sortie}
  Form  (1 hooks)  tests/fixtures/wasted_subtree_render/child_setter.tsx
    warn   wasted-subtree-render  [hook:0]  (line 18:6)  each `change` event in `<Field>` writes
      state `v` and re-renders `<Field>` and `<Heavy>` (at least 3 component renders, including a
      list) although none of their inputs depends on it. […]
       → `<Field>` re-renders (1 component render) although none of its props depends on it (line 18:6)
       → `<Heavy>` re-renders (at least 2 component renders, including a list) although none of its
         props depends on it (line 19:6)
\end{sortie}
(Sortie repliée, conseil final élidé en \og[…]\fg{}.) L'atterrissage est dans \code{Field}, donc
le finding pointe l'élément \code{<Field>} du propriétaire (ligne 18) et le message dit
\og in \code{<Field>}\fg{} (champ \code{place} du \code{Trigger},
\src{src/rules/impls/wasted_subtree_render.rs}{101}{123}). Plus surprenant : \code{<Field>}
lui-même est déclaré gaspillé. C'est exact. Sa seule prop, \code{onChange}, porte
\code{\{Setter(0)\}}, et la question est \code{\{Slot(0)\}} : un setter ne change jamais
d'identité, donc \code{Field} re-rend à chaque frappe pour produire la même sortie.
\end{exemple}

\paragraph{Les effets : listeners et minuteries.} La seconde boucle parcourt la relation
\relation{slot_writers} du composant (\cref{def:slot-writer} ;
\src{src/rules/impls/wasted_subtree_render.rs}{124}{146}). Elle ne garde que les écrivains
\emph{locaux} (\code{w.owner} vide : une écriture à travers le setter d'un autre composant est
comptée chez cet autre), qui écrivent un slot d'état du composant, depuis une région
\code{WriterRegion::Effect(e)}. Pour chacun, elle retient les lignes de la relation
\relation{registrations} (\adr{34}) de l'effet \code{e} dont le callback peut appeler le setter.
Deux prédicats syntaxiques en décident :

\rustsnippet{src/rules/impls/wasted_subtree_render.rs}{315}{333}{peut-il appeler le setter ? derrière un test ?}

\code{may_call} penche vers \og oui\fg{} : un callback qui n'est pas une fonction littérale (un
listener nommé que ce parcours ne suit pas) est supposé écrire. \code{keyed_call} penche vers
\og non gardé\fg{}. Les deux erreurs possibles vont donc dans le même sens : plus de déclencheurs,
et des déclencheurs plus fréquents. Le nom de l'événement est celui de l'enregistrement
(\code{"resize"} dans \code{addEventListener("resize", f)}), à défaut le nom du registrar
(\code{setInterval}).

\begin{exemple}{Une minuterie qui écrit deux slots}{rendu-clock}
Le fichier \file{/tmp/ch21/clock.tsx} :
\begin{lstlisting}[language=TypeScript,caption={Deux écritures dans un même tick},label={lst:rendu-clock}]
function Row({ i }: { i: number }) { return <li>{i}</li>; }
function Agenda() { return <ul>{[1, 2, 3].map((i) => <Row key={i} i={i} />)}</ul>; }

export default function Dashboard() {
  const [now, setNow] = useState(Date.now());
  const [ticks, setTicks] = useState(0);
  useEffect(() => {
    const id = setInterval(() => { setNow(Date.now()); setTicks((t) => t + 1); }, 1000);
    return () => clearInterval(id);
  }, []);
  return (<main><time>{now}</time><Agenda /></main>);
}
\end{lstlisting}
\begin{sortie}
  Dashboard  (3 hooks)  clock.tsx
    warn   wasted-subtree-render  [hook:0]  (line 9:2)  each `setinterval` event writes states `now`
      and `ticks` and re-renders `<Agenda>` (at least 2 component renders, including a list)
      although none of its inputs depends on them. Move `now` and `ticks` and the elements that use
      them into their own component, or build the unrelated elements higher up and pass them in as
      `children`
       → `<Agenda>` re-renders (at least 2 component renders, including a list) although none of its
         props depends on it (line 16:6)
\end{sortie}
Les deux écritures appartiennent au déclencheur \code{Effect(1)} : un seul finding, qui nomme les
deux slots (\og states \code{now} and \code{ticks}\fg{}), positionné sur l'effet (ligne 9).
L'événement s'appelle \code{setinterval} : c'est le nom du registrar, mis en minuscules.
\end{exemple}

Quand un effet vient d'un hook personnalisé greffé dans le composant, le correctif n'appartient
pas au composant : tous les appelants du hook paient le même prix. \code{region_hook}
(\src{src/rules/impls/wasted_subtree_render.rs}{294}{313}) retrouve le nom du hook par la
provenance (\code{hook_provenance}) de l'effet greffé, et le message change de conseil. Sur la
fixture \file{tests/fixtures/wasted_subtree_render/hook_trigger/}, où \code{useWindowSize} écrit
un objet frais à chaque \code{resize} :
\begin{sortie}
  Page  (1 hooks)  tests/fixtures/wasted_subtree_render/hook_trigger/Page.tsx
    warn   wasted-subtree-render  [hook:1]  (tests/fixtures/wasted_subtree_render/hook_trigger/use-size.ts:5:2)
      each `resize` event writes state `size` and re-renders `<List>` (at least 2 component renders,
      including a list) although none of its inputs depends on it. The writes come from
      `useWindowSize`, which does this to every component that calls it: write only when a value its
      callers read changes, or call it from a smaller component
       → `<List>` re-renders (at least 2 component renders, including a list) although none of its
         props depends on it (line 9:48)
\end{sortie}
La position est dans l'autre fichier (\file{use-size.ts}, ligne 5) : c'est celle de l'effet
greffé.

\subsection{Étape 2 : la fréquence, un classement et non une preuve}

Un clic qui ouvre une boîte de dialogue re-rend la page une fois. Une frappe dans un champ la
re-rend à chaque caractère. Par défaut, la règle ne parle que du second cas. La fréquence d'un
événement est calculée par \code{event_frequency} :

\rustsnippet[label={lst:rendu-frequency}]{src/rules/helpers/render_tree.rs}{720}{751}{la fréquence d'un événement sur une cible}

Elle s'appuie sur quatre tables (\src{src/rules/helpers/render_tree.rs}{677}{713}) résumées dans
la \cref{tab:rendu-frequence}. L'arbre de décision est le suivant.
\begin{enumerate}
  \item Un événement de mouvement (\code{mousemove}, \code{scroll}, \code{resize}, \code{drag}\dots)
        ou \code{setinterval} est continu, quelle que soit la cible.
  \item Un événement qui n'est pas de frappe (\code{click}, \code{mouseover}, \code{select}\dots)
        est discret ; un événement clavier (\code{keydown}, \code{keyup}, \code{keypress}) dont
        l'écriture est gardée par un test de l'événement (\code{keyed}) aussi.
  \item Reste un événement de frappe ; la cible tranche. Un \code{<textarea>} est continu ; un
        \code{<input>} l'est si son \code{type} est absent ou textuel (\code{number} et
        \code{range} compris) ; un autre hôte (\code{contenteditable}) l'est pour \code{input} et
        \code{beforeinput}. Un élément de composant l'est si son nom contient un indice
        (\code{Input}, \code{Search}, \code{Editor}, \code{Slider}\dots). Un listener d'effet, sans
        cible, l'est pour \code{input}, \code{keydown} et \code{keyup}.
\end{enumerate}

\begin{table}[htbp]\centering\small
\begin{tabularx}{\linewidth}{>{\raggedright\arraybackslash}p{3.4cm}X}
\toprule
Table & Contenu \\
\midrule
\code{MOTION_EVENTS} & \code{mousemove}, \code{pointermove}, \code{touchmove}, \code{scroll}, \code{wheel}, \code{drag}, \code{dragover}, \code{resize}, \code{selectionchange} \\
\code{TYPING_EVENTS} & \code{change}, \code{input}, \code{keydown}, \code{keyup}, \code{keypress}, \code{beforeinput} \\
\code{TEXT_INPUT_TYPES} & \code{text}, \code{search}, \code{email}, \code{password}, \code{url}, \code{tel}, \code{number}, \code{range} \\
\code{TEXT_COMPONENT_HINTS} & \code{Input}, \code{TextArea}, \code{Textarea}, \code{Search}, \code{Editor}, \code{Slider} \\
\code{KEY_EVENTS} & \code{keydown}, \code{keyup}, \code{keypress} \\
\bottomrule
\end{tabularx}
\caption{Les tables de classement de \file{render_tree.rs}.}
\label{tab:rendu-frequence}
\end{table}

Un événement continu ne suffit pas. Un \code{scroll} qui écrit \code{window.scrollY > 50} écrit à
chaque pas, mais React abandonne toute écriture \code{Object.is}-égale à la valeur courante : le
booléen ne re-rend qu'à ses basculements. La règle exige donc qu'au moins un des slots du
déclencheur ait des valeurs \emph{nombreuses}, au sens du domaine des valeurs :

\rustsnippet{src/domains/impls/state_value.rs}{384}{401}{un slot à peu de valeurs re-rend au rythme de ses transitions}

Un slot est \og à valeurs finies\fg{} s'il ne contient que des booléens, \code{null} ou
\code{undefined}, un ensemble fini de chaînes constantes ou un intervalle entier de largeur au
plus 16, et aucune référence ni setter. C'est la seule lecture du domaine des valeurs que fait
la règle, et elle porte sur le store d'état convergé (\code{result.state_store}). La
fixture \file{scroll_flag.tsx} (un \code{scroll} qui écrit un booléen) est muette par défaut, et
tire avec \code{continuousOnly=false}.

\begin{remarque}
Le texte de \code{keyed} vient du fait \code{gated} de la dépendance de rendu : la part de
l'ensemble qu'une valeur fonction n'atteint que derrière un test de ses propres arguments
(\code{Deps::gated_for}, \src{src/engine/render_deps.rs}{189}{198}). C'est un fait must : un seul
contributeur non gardé retire la source. Sur \file{keyed.tsx}, les trois handlers qui n'écrivent
que sur \code{Enter} ou \code{Escape} sont discrets (muets par défaut) ; seul \code{KeyLog}, qui
écrit \code{e.key} à chaque touche, tire.
\end{remarque}

Le point important est ailleurs : \emph{la fréquence ne prouve rien}. Elle ne décide que de ce qui
est montré par défaut. Classer un événement continu comme discret masque un finding que
\code{continuousOnly=false} rendrait visible ; l'inverse montre un finding par défaut. Aucune des
deux erreurs ne change l'affirmation du finding, qui porte sur les entrées des éléments, pas sur
la fréquence (\adr{41}, conséquences).

\subsection{Étape 3 : les frères gaspillés}

Pour un déclencheur retenu, la question est
$\mathit{rel} = \{\code{Slot}(l)\}_l \cup \{\code{Module}(x)\}_x$ : les slots du lot, et les noms
de module que le lot écrit à côté d'eux (un handler qui fait \code{cache.x = v; setV(v)} change
aussi ce que lisent les lecteurs de \code{cache}). \code{wasted_siblings} y répond :

\rustsnippet[label={lst:rendu-wasted-siblings}]{src/rules/helpers/render_tree.rs}{266}{326}{les éléments que le lot ne peut pas changer}

On la lit en quatre temps.
\begin{enumerate}
  \item \textbf{Touché.} Un site est touché si sa garde, l'une de ses props ou son spread dépend
        de la question. Pour un provider de contexte (\code{provides}), seule la garde compte :
        sa \code{value} ne va pas à ses enfants par les props, elle est suivie autrement.
  \item \textbf{Atteint.} Un site est écarté s'il est touché \emph{ou} si l'un des éléments dans
        lesquels il est imbriqué (la chaîne des \code{parent}) l'est. Un
        \code{<Panel value=\{text\}><Heavy /></Panel>} ne rend pas \code{<Heavy>} candidat :
        \code{Panel} pourrait lui transmettre la valeur, par contexte ou par \code{cloneElement}.
  \item \textbf{Résolu.} Seul un site dont \code{resolve} trouve le composant
        (\cref{lst:rendu-resolve}) est candidat. Un élément opaque n'est jamais déclaré gaspillé.
  \item \textbf{Non utilisateur.} La question transportée (\code{carried}) rassemble les
        contextes qui portent la valeur jusqu'à ce site (\code{contexts_at}) et les noms de
        module écrits. Si l'enfant \emph{utilise} l'un d'eux, ou s'il n'a pas de résumé, il est
        écarté : un consommateur du contexte re-rend de toute façon, à bon droit.
\end{enumerate}

\begin{figure}[htbp]\centering
\begin{tikzpicture}[node distance=5mm and 7mm,
    q/.style={bloc,minimum width=30mm,align=center},
    refus/.style={draw=ra-gray,rounded corners,align=center,font=\small,inner sep=3pt}]
  \node[q] (t) {site $i$ de $o$};
  \node[q,below=of t] (r) {touché, ou un \code{parent}\\ touché ?};
  \node[q,below=of r] (s) {\code{resolve} trouve\\ un composant ?};
  \node[q,below=of s] (u) {l'enfant utilise\\ \code{carried} ?};
  \node[q,below=of u,draw=ra-blue,very thick] (w) {gaspillé :\\ \code{subtree_renders}};
  \node[refus,right=12mm of r] (r1) {écarté : l'élément\\ dépend du lot};
  \node[refus,right=12mm of s] (s1) {écarté : inconnu\\ = usage};
  \node[refus,right=12mm of u] (u1) {écarté : consommateur\\ (ou pas de résumé)};
  \draw[fleche] (t) -- (r);
  \draw[fleche] (r) -- node[etiquette,left] {non} (s);
  \draw[fleche] (s) -- node[etiquette,left] {oui} (u);
  \draw[fleche] (u) -- node[etiquette,left] {non} (w);
  \draw[fleche] (r) -- node[etiquette,above] {oui} (r1);
  \draw[fleche] (s) -- node[etiquette,above] {non} (s1);
  \draw[fleche] (u) -- node[etiquette,above] {oui} (u1);
\end{tikzpicture}
\caption{Le filtre de \code{wasted_siblings} sur un site du propriétaire. Toute sortie à droite
est un refus de conclure ; seule la sortie du bas produit une affirmation.}
\label{fig:rendu-filtre}
\end{figure}

La \cref{fig:rendu-filtre} résume ce filtre. Chaque question y a deux issues, et l'issue
prudente est toujours celle qui écarte le site. C'est l'\cref{inv:tout-inconnu-usage}, écrit en
code.

\subsection{Étape 4 : compter les rendus}

Pour un candidat, \code{subtree_renders} compte les rendus de composant qu'il coûte :

\rustsnippet[label={lst:rendu-subtree}]{src/rules/helpers/render_tree.rs}{328}{372}{le coût d'un sous-arbre gaspillé}

Le candidat compte pour un. Chaque élément de composant qu'il construit ajoute le coût de son
propre sous-arbre, sauf s'il consomme ce que le lot change (il re-rend à bon droit, et ce qui est
sous lui aussi) et sauf s'il est un provider (qui n'est pas un rendu de composant). Un site
\code{in_list} ne compte \emph{qu'une fois} et allume le drapeau \code{list} : le vrai nombre
d'instances est inconnu, peut-être grand, peut-être nul. La récursion et la profondeur
(\code{MAX_DEPTH}) arrêtent la descente en comptant le composant pour un. Le compte est donc
présenté comme une borne inférieure (\og at least N component renders, including a list\fg{}),
avec une exception discutée à la \cref{sec:regles-rendu-rsc-limites}.

\subsection{Étape 5 : seuil et message}

Le seuil est appliqué en \src{src/rules/impls/wasted_subtree_render.rs}{185}{190} : le
déclencheur produit un finding si au moins un élément est gaspillé et si le total atteint
\code{minWastedRenders}, \emph{ou} si l'un des sous-arbres contient une liste. Le message assemble le déclencheur
(premier événement continu, sinon le premier), les éléments gaspillés, le coût, puis l'un des deux
correctifs structurels : déplacer l'état et ses usagers dans leur propre composant, ou construire
les éléments indépendants plus haut et les passer en \code{children}. Chaque élément gaspillé
reçoit une note \code{Step::Rerender}, visible avec \code{--trace}. Le finding est porté par le
label du premier slot et par le span du handler (ou de l'élément \code{via}, ou de l'effet).

\section{Des options sur une règle native}
\label{sec:regles-rendu-rsc-options}

Jusqu'à \adr{41}, aucune règle native n'acceptait de paramètre (\adr{22} §4). Une règle
classée par coût change la donne : savoir à partir de combien de rendus inutiles un sous-arbre
mérite d'être signalé, ou si un clic compte, est un choix d'équipe, pas un fait que l'analyseur
peut établir. La règle déclare donc deux options typées :

\rustsnippet{src/rules/impls/wasted_subtree_render.rs}{33}{48}{les deux options de la règle}

Le type \code{OptionSpec} (\src{src/rules/api/query.rs}{272}{288}) ne connaît que deux formes,
un entier non signé borné et un booléen. Le registre valide chaque valeur configurée
\emph{avant} l'analyse, et bruyamment : une clé inconnue ou une valeur hors bornes est une erreur
d'usage, jamais un réglage ignoré. Les accesseurs \code{RuleConfig::uint} et \code{flag}
(\src{src/rules/api/query.rs}{247}{269}) peuvent alors rendre la valeur ou le défaut sans
revérifier.

\begin{shell}
reactant check typing.tsx --rule wasted-subtree-render \
  --rule-option wasted-subtree-render:minWastedRenders=0
\end{shell}
\begin{sortie}
[error] option `minWastedRenders` of rule `wasted-subtree-render` must be an integer between 1 and 1000
\end{sortie}
Le code de sortie est 2 (erreur d'usage). Les options se donnent aussi dans le fichier de
configuration ; \code{--rule-option} l'emporte sur la valeur de même clé (\adr{41} §4). Avec,
dans \file{/tmp/ch21/cfg/reactant.config.json} :
\begin{lstlisting}[language=JSON]
{ "rules": { "wasted-subtree-render":
    { "options": { "continuousOnly": false, "minWastedRenders": 1 } } } }
\end{lstlisting}
le composant \code{HoverCard} suivant, muet par défaut, produit deux findings :
\begin{lstlisting}[language=TypeScript,caption={Un survol à côté d'un sous-arbre (\file{/tmp/ch21/cfg/hover_inline.tsx})},label={lst:rendu-hover}]
function Heavy() { return <section>{[0, 1, 2].map((i) => <span key={i}>{i}</span>)}</section>; }
export default function HoverCard() {
  const [hover, setHover] = useState(false);
  return (
    <div className={hover ? "on" : "off"}
         onMouseOver={() => setHover(true)} onMouseOut={() => setHover(false)}>
      <Heavy />
    </div>
  );
}
\end{lstlisting}
\begin{sortie}
  HoverCard  (3 hooks)  hover_inline.tsx
    warn   wasted-subtree-render  [hook:0]  (line 8:9)  each `mouseover` event writes state `hover`
      and re-renders `<Heavy>` (1 component renders) although none of its inputs depends on it. […]
    warn   wasted-subtree-render  [hook:0]  (line 8:44)  each `mouseout` event writes state `hover`
      and re-renders `<Heavy>` (1 component renders) although none of its inputs depends on it. […]
\end{sortie}
Il fallait les deux options. \code{mouseover} n'est pas un événement continu et \code{hover} est
un booléen : \code{continuousOnly=false} est nécessaire. Le \code{.map} de \code{Heavy} ne
construit que des éléments hôtes (\code{<span>}), donc aucun site \code{in_list}, et le coût est
d'un seul rendu : \code{minWastedRenders=1} est nécessaire aussi. Les deux handlers sont deux
déclencheurs distincts, donc deux findings.

\begin{decision}[ADR-041 §4]
\textbf{Décision.} Une règle native peut déclarer des options typées (\code{Rule::options}), que
le registre valide ; \code{--rule-option <rule>:<key>=<value>} bat la configuration. Les défauts
(2 et \code{true}) viennent d'un balayage du corpus.
\textbf{Ce qui est amendé.} \adr{22} §4 (\og natives declare no params in v1\fg{}) : les seuils
d'une règle classée par coût sont l'affaire de l'équipe. Une règle native qui ne déclare rien
continue de refuser toute option.
\end{decision}

\begin{piege}
Une option ne peut qu'ajuster \emph{ce qui est montré}, jamais ce qui est affirmé. Baisser
\code{minWastedRenders} ou couper \code{continuousOnly} fait apparaître des findings dont
l'affirmation (\og aucune entrée ne dépend de l'état\fg{}) est calculée de la même façon. Aucune
option ne relâche l'\cref{inv:tout-inconnu-usage}.
\end{piege}

\section{Cas limites : ce que les fixtures épinglent}
\label{sec:regles-rendu-rsc-cas}

Les fixtures de \file{tests/fixtures/wasted_subtree_render/} sont le contrat de la règle. Cinq
d'entre elles (\file{typing}, \file{composer}, \file{dropzone}, \file{click}, \file{children})
viennent de \file{scripts/rerender-bench/}, où leur nombre de rendus a été mesuré sous React~19
(en-tête de \file{tests/wasted_subtree_render.rs}). On en retient quatre familles.

\paragraph{Le correctif \og children\fg{}.} C'est le cas limite le plus instructif, parce qu'il
illustre le premier point d'arrêt de la cascade (encadré du début du chapitre) :

\tsxsnippet[label={lst:rendu-children}]{tests/fixtures/wasted_subtree_render/children.tsx}{3}{16}{l'élément lourd est construit plus haut}

\code{<Heavy />} est construit dans le rendu d'\code{App}, qui ne se ré-exécute pas quand
\code{hover} change ; \code{HoverCard} reçoit l'objet élément par sa prop \code{children} et le
renvoie tel quel, et React saute son rendu. Côté analyse, c'est immédiat : le résumé de
\code{HoverCard} n'a \emph{aucun} site (l'élément est un site d'\code{App}), donc
\code{wasted_siblings} n'a rien à examiner. La fixture est muette même avec
\code{continuousOnly=false} et \code{minWastedRenders=1} (sortie observée : \og ✓ 1 file(s) no issues found.\fg{}), alors que
la variante où \code{HoverCard} construit lui-même \code{<Heavy />}
(\cref{lst:rendu-hover}) tire. La dépendance de rendu distingue ainsi, sans règle spéciale,
\og construire\fg{} et \og transmettre\fg{} un élément.

\paragraph{Le provider.} Dans \file{provider.tsx}, \code{App} rend
\code{<Ctx.Provider value=\{text\}>} autour de \code{<Panel><Reader /></Panel>}, et
\code{Reader} lit le contexte. Le provider n'est pas touché (seule sa garde compterait), mais
\code{contexts_at} ajoute \code{Context(Ctx)} à la question transportée pour tout site imbriqué
dans lui. \code{<Reader>}, site d'\code{App} imbriqué dans \code{<Panel>}, lui-même imbriqué dans
le provider, utilise ce contexte : il est écarté, rien n'est affirmé sur lui. \code{<Panel>}, en
revanche, ne lit rien : il \emph{est} gaspillé, pour un rendu. La fixture n'est muette par défaut
qu'à cause du seuil ; avec \code{minWastedRenders=1} et \code{continuousOnly=false}, on observe
\og re-renders \code{<Panel>} (1 component renders)\fg{}. Le nom du test qui l'épingle,
\code{children_from_above_and_provider_subtrees_are_not_claimed}, décrit donc le sort de
\code{<Reader>}, pas celui de \code{<Panel>}. La fixture
\file{context.tsx} montre l'autre face (\adr{41}, \issue{145}) : les éléments qu'un provider
enveloppe \emph{peuvent} être gaspillés quand ils ne consomment pas le contexte, et le test
\code{a_provider_is_followed_to_its_consumers} vérifie qu'un consommateur, direct ou à travers un
hook, n'est jamais compté.

\paragraph{La garde.} Dans \file{click.tsx}, \code{\{open \&\& <Dialog onClose=\{…\} />\}} a une
garde \code{\{Slot(0)\}} : l'élément existe ou non selon l'état, il est touché, jamais candidat.
Seul \code{<BigList>} est déclaré gaspillé, et seulement avec \code{continuousOnly=false} :
\begin{sortie}
  Page  (3 hooks)  tests/fixtures/wasted_subtree_render/click.tsx
    warn   wasted-subtree-render  [hook:0]  (line 14:24)  each `click` event writes state `open` and
      re-renders `<BigList>` (at least 2 component renders, including a list) […]
    warn   wasted-subtree-render  [hook:0]  (line 15:15)  each `click` event in `<Dialog>` writes
      state `open` and re-renders `<BigList>` (at least 2 component renders, including a list) […]
\end{sortie}
Deux déclencheurs : le bouton d'ouverture du propriétaire, et le bouton de fermeture de
\code{Dialog}, où le setter a atterri à travers la closure \code{onClose}.

\paragraph{L'élément opaque.} Enveloppons \code{Sidebar} de la \cref{lst:rendu-recherche} dans
\code{memo} (\file{/tmp/ch21/memo.tsx}) : la règle devient muette. Ce n'est pas qu'elle
reconnaisse la barrière \code{memo} ; c'est que \code{memo(function Sidebar() \{…\})} ne se résout
vers aucun composant enregistré (\issue{64}), et qu'un élément opaque n'est jamais candidat. Le
résultat est juste ici, mais pour une raison qui vaut aussi pour une bibliothèque \emph{sans}
\code{memo} : un sous-arbre lourd importé d'un paquet, à côté d'un champ rapide, n'est pas signalé
(\file{docs/limitations.md}, \og Render cascades stop at what the analysis cannot see
into\fg{}).

\begin{table}[htbp]\centering\small
\begin{tabularx}{\linewidth}{>{\raggedright\arraybackslash}p{3.2cm}X>{\raggedright\arraybackslash}p{3.3cm}}
\toprule
Fixture & Ce qu'elle épingle & Verdict par défaut \\
\midrule
\file{typing.tsx} & frappe à côté d'un sous-arbre avec liste & \Warning{} \\
\file{typing_fixed.tsx} & état déplacé dans \code{<Editor>} & muet \\
\file{child_setter.tsx} & setter passé à un enfant (\issue{146}) & \Warning{}, \code{<Field>} compris \\
\file{keyed.tsx} & écriture derrière \code{e.key} : discret (\issue{148}) & \Warning{} sur \code{KeyLog} seul \\
\file{click.tsx} & clic discret, garde & muet ; 2 \Warning{} sans \code{continuousOnly} \\
\file{scroll_flag.tsx} & booléen écrit par \code{scroll} & muet ; \Warning{} sans \code{continuousOnly} \\
\file{children.tsx} & élément reçu en \code{children} & muet dans tous les cas \\
\file{provider.tsx} & consommateur sous un provider & muet (seuil : \code{<Panel>} coûte 1) \\
\file{context.tsx} & provider suivi jusqu'aux consommateurs (\issue{145}) & 2 \Warning{} \\
\file{hook_trigger/} & effet venu d'un hook personnalisé & \Warning{} nommant le hook \\
\file{module_write.tsx} & nom de module écrit à côté de l'état (\issue{147}) & 3 \Warning{} ; 4 sans \code{continuousOnly} \\
\bottomrule
\end{tabularx}
\caption{Quelques fixtures de la règle et leur verdict, rejoués au commit photographié.}
\label{tab:rendu-fixtures}
\end{table}
% Note : composer, dropzone, pointer, spread_input, wrapped_input, checkbox ne figurent pas dans la table (non rejoués).

\section{Niveau et soundness d'une affirmation d'absence}
\label{sec:regles-rendu-rsc-soundness}

\paragraph{Ce que dit un finding.} Un finding de \regle{wasted-subtree-render} est la conjonction
de trois affirmations : (a) le déclencheur écrit les slots nommés ; (b) chaque écriture re-rend le
propriétaire, donc chaque élément gaspillé ; (c) aucune entrée de ces éléments (props, spread,
garde, éléments englobants, contextes portés, noms de module écrits) ne dépend de ce que le lot
change. La première vient des relations du moteur ; la deuxième est la sémantique de React
(\cref{chap:react-modele}) ; la troisième est l'affirmation d'absence.

\paragraph{Polarité renversée.} Pour une affirmation d'absence, sur-approximer les
\emph{dépendances} est la direction prudente : une dépendance de trop fait taire la règle (un
finding de moins), une dépendance oubliée fabriquerait une affirmation fausse. C'est exactement ce
que garantit l'analyse avant de \file{render_deps.rs} : un may-ensemble, joint par union, dégradé
en top au plafond, jamais en plus petit. Et c'est ce que prolonge l'\cref{inv:tout-inconnu-usage}
au niveau de l'arbre : tout ce qui n'est pas vu compte comme un usage.

\begin{theoreme}{Correction de l'affirmation d'absence}{rendu-absence}
Supposons que, pour chaque composant $C$, le résumé \code{render_deps}$(C)$ sur-approxime les
sources dont dépendent la sortie, les effets et les éléments de $C$ (hypothèse de
\cref{def:render-dep}), et qu'aucun appel opaque ne lise ni n'écrive un binding de module
(\adr{41} §2). Alors tout élément $e$ rapporté par \code{wasted_siblings}$(o, \mathit{rel})$ a
des props, un spread et une garde indépendants de $\mathit{rel}$, n'est imbriqué dans aucun élément
qui en dépend, se résout vers un composant $C$ qui n'utilise pas ce qui est porté jusqu'à lui, et
tout sous-arbre qu'il compte est fait d'éléments résolus non consommateurs.
\end{theoreme}

\noindent\textit{Esquisse.} Chaque exclusion de la \cref{fig:rendu-filtre} teste \code{touches}
sur un ensemble qui contient les vraies dépendances ; un test négatif sur un sur-ensemble est
négatif sur l'ensemble. Les cas où l'ensemble n'est pas connu (non résolu, top, récursion,
profondeur, contexte innommable) sont tous rangés du côté \og usage\fg{}.

\paragraph{Ce qui fuit.} Deux hypothèses restent, toutes deux écrites dans \adr{41}. La première
est l'appel opaque : une fonction que l'analyse ne voit pas est supposée ne lire ni écrire aucun
binding de module (\issue{51}, \issue{52}). La seconde est la correspondance par orthographe des
noms de module entre fichiers : une collision donne \emph{moins} de findings, jamais un faux. Si
la première est violée, la règle peut affirmer à tort qu'un élément est gaspillé : c'est un faux
positif, toléré au niveau \Warning{}, et c'est la seule direction où elle peut se tromper sur
l'affirmation.

\paragraph{Et les faux négatifs ?} Le prix de la prudence est une sous-déclaration : un sous-arbre
\code{memo} ou de bibliothèque, un élément de nom ambigu, une cascade à travers un contexte
importé d'un paquet ne sont jamais signalés. Le projet les classe comme limites de précision
(\issue{64} porte l'étiquette \code{precision-fn}, \og sound but incomplete\fg{}) et non comme
bogues de soundness : l'invariant \og faux négatifs interdits\fg{} (\cref{def:soundness}) porte
sur les défauts que l'analyse \emph{prétend} exclure, et cette règle ne prétend rien sur un
élément qu'elle n'a pas pu résoudre.

\begin{decision}[ADR-041 §1–§3]
\textbf{Décision.} La dépendance de rendu est une analyse avant séparée, pas un champ de
\code{StateValue}. \textbf{Alternatives refusées.} (1) Réutiliser \code{Versioned(S)}
(\adr{17}) : ce fait ne vit que sur la composante \code{reference}, un booléen ou
\code{text.length} n'y portent rien, et une allocation est \code{PerRender} quelle que soit son
origine. (2) Ajouter la dépendance comme composante du produit : toutes les comparaisons de
valeurs existantes en seraient changées (\regle{unnecessary-rerender} compare \code{arg == init},
la clé du \code{ComponentCache} compare les props). \textbf{Niveau.} \Warning{} : les rendus
supplémentaires sont certains, leur coût ne l'est pas (définition du \Warning{} dans
\file{CLAUDE.md}).
\end{decision}

\begin{remarque}
Pourquoi pas \Error{} quand le sous-arbre est énorme ? Parce que le niveau \Error{} exige un
\emph{défaut} certain (\cref{def:certified}). Un rendu inutile n'est pas un défaut : c'est un
coût, et la règle n'a aucun moyen de le chiffrer (un composant peut être trivial à rendre). La
règle n'a d'ailleurs pas de \code{safe_check} : une assurance \og aucun rendu gaspillé\fg{}
reposerait sur la même borne inférieure, et n'aurait pas de sens pour un élément non résolu.
\end{remarque}

% ===========================================================================
\section{Server Components : le graphe serveur de Next.js}
\label{sec:regles-rendu-rsc-graphe}

\begin{react}[Server Components et \code{"use client"}]
Sous l'App Router de Next.js, les fichiers réservés d'un dossier \file{app/} (\file{page.tsx},
\file{layout.tsx}\dots) sont des \emph{Server Components} : React les rend une fois, côté
serveur, et envoie le résultat. Il n'y a ni état, ni commit, ni effets ; React n'a aucun hook à
leur fournir, et un \code{useState} y fait échouer le rendu. Une directive \code{"use client"},
écrite dans le \emph{prologue} d'un module (avant toute instruction), ouvre une frontière :
ce module, et tout ce qu'il importe, est compilé pour le client. La propriété \og être serveur\fg{}
n'est donc pas celle d'un fichier isolé, mais de \emph{la façon dont on l'atteint} par les
imports. Un module importé des deux côtés est compilé deux fois, et sa copie serveur doit
pouvoir s'exécuter. Voir \cref{chap:react-modele}.
\end{react}

\regle{server-component-hook} répond à une question que le domaine abstrait ignore
complètement : dans quel environnement ce module s'exécute-t-il ? Elle s'appuie sur deux faits
de nature différente. Le premier est une \emph{convention de nommage} : quels fichiers sont des
entrées serveur. Le second est un \emph{graphe} : les arêtes d'import résolues, que le front-end
enregistre dans les \code{ModuleFacts} de chaque fichier (\cref{chap:frontend-detection},
\cref{lst:frontend-modulefacts}) : le prologue de directives, verbatim, et les fichiers importés
que le résolveur a pu associer à un vrai fichier, \code{import type} exclu.

\subsection{Les entrées : une convention de noms}

\rustsnippet{src/project/nextjs.rs}{13}{38}{les entrées App Router}

Il faut les deux moitiés de la convention : un nom réservé \emph{et} un ancêtre nommé
\file{app}. Le nom seul désignerait n'importe quel \file{components/page.tsx}. \code{error} et
\code{global-error} sont absents à dessein : Next.js exige qu'ils soient des Client Components.
Les groupes de routes entre parenthèses passent sans traitement particulier : le test
\code{entry_kinds_need_stem_and_app_dir} accepte \file{/r/src/app/(marketing)/x/layout.tsx}.

\subsection{Le graphe serveur : une atteignabilité avec frontière}

\rustsnippet[label={lst:rendu-server-modules}]{src/project/nextjs.rs}{40}{55}{les modules du graphe serveur}

\code{server_modules} part de toutes les entrées (les \code{seeds} du parcours), puis appelle le mécanisme unique
de la table, \code{ModuleTable::reachable_from(seeds, Some("use client"))}
(\cref{lst:frontend-reachable}, \src{src/ir/module.rs}{91}{126}). C'est un parcours en largeur
sur les arêtes d'import ; un module qui déclare la frontière n'est ni visité ni traversé, qu'il
soit un point de départ ou non. Une page qui porte elle-même \code{"use client"} n'est donc pas serveur,
et un module que seul un module client importe ne l'est pas non plus. Le coût est
$O(V + E)$ sur le graphe de modules.

\begin{definition}{Module serveur}{module-serveur}
Un module $m$ est \terme{module serveur} s'il existe un chemin d'arêtes d'import résolues
$e = m_0 \to m_1 \to \dots \to m_k = m$ tel que $e$ est une entrée App Router
(\code{server_entry_kind}$(e) \neq$ \code{None}) et qu'aucun $m_i$ ($0 \le i \le k$) ne déclare
\code{"use client"} dans son prologue.
\end{definition}

\begin{exemple}{Un module partagé et une frontière}{rsc-partage}
Le projet \file{/tmp/ch21/nx/} (un \file{next.config.js}, un \file{tsconfig.json}) contient
cinq modules :
\begin{lstlisting}[language=TypeScript,caption={Le projet \file{/tmp/ch21/nx/} (fichiers concaténés)},label={lst:rsc-nx}]
// app/page.tsx
import { Search } from "../components/search";
import { Price } from "../components/price";
export default function Home() { return <main><Search /><Price value={3} /></main>; }

// app/shop/page.tsx
import { useSession } from "next-auth/react";
import { useFormat } from "../../lib/format";
export default function Shop() {
  const session = useSession();
  const total = useFormat(10);
  return <p>{total}{String(session)}</p>;
}

// components/search.tsx
"use client";
import { useState } from "react";
import { Price } from "./price";
export function Search() {
  const [q, setQ] = useState("");
  return <div><input value={q} onChange={(e) => setQ(e.target.value)} /><Price value={q.length} /></div>;
}

// components/price.tsx
import { useId } from "react";
import { useFormat } from "../lib/format";
export function Price({ value }: { value: number }) {
  const id = useId();
  const text = useFormat(value);
  return <span id={id}>{text}</span>;
}

// lib/format.ts
export function useFormat(n: number) { return n.toFixed(2); }
\end{lstlisting}
La \cref{fig:rsc-graphe} montre le graphe et la frontière. Les points de départ sont les
deux \file{page.tsx}. Depuis \file{app/page.tsx}, le parcours atteint \file{price.tsx}, mais s'arrête
à \file{search.tsx}, qui déclare la frontière. \file{price.tsx} est importé des deux côtés : il
est serveur (par \file{app/page.tsx}), quoique \file{search.tsx} l'importe aussi.
\end{exemple}

\begin{figure}[htbp]\centering
\begin{tikzpicture}[node distance=9mm and 14mm,
    mod/.style={bloc,minimum width=30mm,font=\small},
    seed/.style={mod,very thick,draw=ra-blue},
    cli/.style={mod,fill=white}]
  \node[seed] (home) {\file{app/page.tsx}\\[-2pt]{\scriptsize entrée \code{page}}};
  \node[seed,right=36mm of home] (shop) {\file{app/shop/page.tsx}\\[-2pt]{\scriptsize entrée \code{page}}};
  \node[cli,below left=14mm and -4mm of home] (search) {\file{components/search.tsx}\\[-2pt]{\scriptsize\code{"use client"}}};
  \node[mod,below right=14mm and -4mm of home] (price) {\file{components/price.tsx}\\[-2pt]{\scriptsize serveur (partagé)}};
  \node[mod,below=14mm of shop] (fmt) {\file{lib/format.ts}\\[-2pt]{\scriptsize serveur}};
  \draw[fleche] (home) -- (price);
  \draw[fleche] (home) -- (search);
  \draw[fleche] (shop) -- (fmt);
  \draw[fleche] (price) -- (fmt);
  \draw[flechemay] (search) -- (price);
  \begin{scope}[on background layer]
    \node[draw=ra-amber,thick,dashed,rounded corners,inner sep=3mm,fit=(search),
          label={[etiquette,text=ra-amber]below:frontière client}] {};
    \node[fill=ra-blue!8,rounded corners,inner sep=2.5mm,fit=(home)(shop)(price)(fmt)] {};
  \end{scope}
\end{tikzpicture}
\caption{Le graphe de modules de \file{/tmp/ch21/nx/}. Fond bleu : \code{server_modules}. La
frontière \code{"use client"} (pointillés orangés) arrête le parcours : \file{search.tsx} n'est
pas visité. L'arête pointillée de \file{search.tsx} vers \file{price.tsx} n'est jamais suivie
depuis le serveur ; \file{price.tsx} est serveur par l'arête directe de la page.}
\label{fig:rsc-graphe}
\end{figure}

\begin{piege}
Une directive n'en est une que dans le prologue. Dans \file{/tmp/ch21/nx3/components/toggle.tsx},
\begin{lstlisting}[language=TypeScript]
import { useState as useLocal } from "react";
"use client";
export function Toggle() { const [on, setOn] = useLocal(false); /* … */ }
\end{lstlisting}
la chaîne \code{"use client"} suit une instruction d'import : oxc ne la range pas dans
\code{Program::directives}, et Next.js ne la lit pas davantage. Le module est serveur (il est
importé par \file{app/page.tsx}) et la règle tire, en nommant le hook par son origine malgré
l'alias :
\begin{sortie}
  Toggle  (2 hooks)  nx3/components/toggle.tsx
    warn   server-component-hook  [hook:0]  (line 5:8)  `useState` is called in a Server Component.
      this module is imported into the App Router's server graph and no `"use client"` directive
      covers it, […]
\end{sortie}
\end{piege}

\section{L'algorithme de server-component-hook}
\label{sec:regles-rendu-rsc-sch}

La règle tient en trois portes et un message (\src{src/rules/impls/server_component_hook.rs}{129}{193}).

\paragraph{Porte 1 : le programme utilise-t-il la directive ?} Sans un seul \code{"use client"}
dans les modules abaissés (\code{ModuleTable::any_declares}), la règle se tait. Un code qui
n'écrit jamais la directive n'est pas un code Server Components, et qualifier ses modules de
\og serveur\fg{} serait une fiction (\src{src/ir/module.rs}{82}{89}). Le projet
\file{/tmp/ch21/nx2/} est la copie de \file{nx/} sans la directive de \file{search.tsx} : aucun
finding (\og ✓ 6 file(s) no issues found.\fg{}), alors que sous Next.js tous ses hooks seraient
fautifs.

\paragraph{Porte 2 : le fichier du composant est-il serveur ?} \code{server_modules(table)} doit
contenir \code{comp.file}. Dans \file{nx/}, \code{Search} est client, il n'est pas examiné.

\paragraph{Porte 3 : quels hooks sont des preuves ?} Chaque \code{HookEntry} du composant passe
par \code{client_only_hook}, avec sa ligne de provenance (\code{HookProvenance}) :

\rustsnippet[label={lst:rsc-client-only}]{src/rules/impls/server_component_hook.rs}{67}{111}{qu'est-ce qui prouve un hook client ?}

Les cinq hooks modélisés (\code{State}, \code{Effect}, \code{Memo}, \code{Callback}, \code{Ref})
sont des preuves par leur \emph{genre} : un \code{useReducer}, rangé en \code{State}, en est une,
et il est nommé par son origine (\code{origin_hook}), pas par l'alias local ni par le genre. Un
hook \code{Custom} n'est une preuve que si son origine figure dans l'une de deux listes fermées :

\rustsnippet{src/rules/impls/server_component_hook.rs}{10}{47}{les hooks documentés comme client}

La liste des paquets est indexée par le spécificateur d'import : un \code{useRouter} du projet ne
prouve rien, celui de \code{next/navigation} si (test \code{the_package_scope_is_load_bearing}).
Un \code{HookEntry::Handler} (un \code{onClick} sur un élément hôte) n'est pas un hook ; c'est une
autre erreur de Next.js, que la règle ne traite pas.

\begin{exemple}{Ce qui n'est pas une preuve}{rsc-preuves}
Reprenons \file{/tmp/ch21/nx/} (\cref{lst:rsc-nx}) :
\begin{sortie}
  Price  (2 hooks)  nx/components/price.tsx
    warn   server-component-hook  [hook:0]  (line 5:8)  `useId` is called in a Server Component. this
      module is imported into the App Router's server graph and no `"use client"` directive covers
      it, so React renders it on the server, where hooks do not exist; add `"use client"` at the
      top of the file, or move the stateful part into a child component that declares it
  Search  (2 hooks)  nx/components/search.tsx  ✓
  Shop  (2 hooks)  nx/app/shop/page.tsx  ✓
\end{sortie}
(Commande : \code{check nx --rule server-component-hook --show-clean} ; les lignes
\code{suspended analysis-limit} de \code{--info} sont omises.) \code{useId} est dans
\code{REACT_CLIENT_HOOKS} : \code{Price}, module partagé, est signalé. \code{Shop} ne l'est pas :
\code{useSession} vient de \code{next-auth/react}, paquet absent de la liste, et
\code{useFormat} est un utilitaire au nom en \code{use} qui n'appelle aucun hook. Le premier cas
est un faux négatif réel (\code{useSession} est un hook client) ; le second est exactement le
faux positif que la liste fermée évite.
\end{exemple}

\paragraph{Le message.} Les noms sont dédoublonnés dans l'ordre de première apparition : trois
\code{useState} sont une seule directive manquante. Il y a \emph{un} finding par composant, qui
nomme au plus trois hooks puis \og and N more\fg{}, et qui distingue une entrée
(\og this file is an App Router \code{page}\fg{}) d'un module atteint transitivement. Sur la
fixture \file{tests/fixtures/next_project/} :
\begin{sortie}
  HomePage  (2 hooks)  tests/fixtures/next_project/src/app/page.tsx
    warn   server-component-hook  [hook:0]  (line 7:8)  `useState` is called in a Server Component.
      this file is an App Router `page` and no `"use client"` directive covers it, […]
  Sidebar  (2 hooks)  tests/fixtures/next_project/src/components/sidebar.tsx
    warn   server-component-hook  [hook:0]  (line 8:8)  `usePathname`, `useState`, `useMemo` are
      called in a Server Component. this module is imported into the App Router's server graph
      and no `"use client"` directive covers it, […]
\end{sortie}
\code{Sidebar} n'a que deux hooks à la surface (\code{usePathname}, \code{useNavItems}) ;
\code{useState} et \code{useMemo} viennent de \code{useNavItems}, greffé par inlining. Le
\code{Counter} client de la même fixture garde son finding \regle{infinite-loop} : les autres
règles ne sont pas supprimées dans un projet Next (test \code{a_client_component_is_not_reported}).

\paragraph{L'assurance.} \code{safe_check} (\src{src/rules/impls/server_component_hook.rs}{118}{127})
ne publie \og this Server Component calls no client-only hook\fg{} que pour un composant qui passe
les portes 1 et 2. Sur un composant client, l'assurance serait vide : il n'y avait rien à
vérifier.

\section{Pourquoi Warning et pas Error}
\label{sec:regles-rendu-rsc-niveau}

Un \code{useState} dans une page sans directive fait échouer le rendu à coup sûr. Pourquoi la
règle ne produit-elle pas une \Error{} ? Parce qu'une \Error{} doit être prouvée dans le domaine,
par un jeton \code{Certified} que seules les primitives must savent frapper
(\cref{def:certified}), et que les deux faits dont dépend la conclusion n'appartiennent pas au
domaine.
\begin{itemize}
  \item \textbf{L'ensemble des entrées} vient d'une convention de noms de fichiers. Un projet peut
        être configuré autrement ; un dossier \file{app/} peut ne pas être celui du routeur.
  \item \textbf{Le graphe} vient des arêtes que le résolveur a trouvées. Une arête manquante
        (alias non résolu) laisse son sous-graphe non classé ; une arête de trop (un import
        \code{\{ type T \}} au niveau du spécificateur, compté comme arête de valeur,
        \cref{chap:frontend-detection}) pourrait rattacher au serveur un module qui n'y est pas.
\end{itemize}
Le projet appelle \og incertain\fg{} un fait qui ne repose pas sur le domaine ; c'est la
définition du \Warning{}. On remarquera que les deux règles du chapitre plafonnent au même niveau
pour des raisons opposées : \regle{wasted-subtree-render} a la preuve mais pas le défaut (un coût
n'est pas un bogue) ; \regle{server-component-hook} a le défaut mais pas la preuve.

\begin{decision}[ADR-026 §4]
\textbf{Décision.} Les modules serveur sont \emph{analysés}, comme tout composant. La règle est un
\Warning{}, gardée par l'usage de la directive, alimentée par les seuls hooks prouvés client, un
finding par composant.
\textbf{Alternatives refusées.} (1) \emph{Sauter} les Server Components : la qualité de serveur
dépend d'arêtes peut-être manquantes, et sauter un module à tort efface ses bogues, un faux négatif.
Les analyser ne coûte rien en soundness : sans état, un Server Component est sur-approximé comme
tout composant aux props top. (2) Une primitive must qui frapperait un \code{Certified} :
\og it would dress a path convention as a domain proof\fg{}. (3) Supprimer les findings des autres
règles dans un module serveur : toute erreur de classement deviendrait un faux négatif. (4)
Accepter tout \code{useX} comme preuve : nombre d'utilitaires nommés \code{use…} n'appellent
aucun hook.
\textbf{Résultat.} Zéro finding sur les corpus Next.js (qui compilent), et un finding correct sous
mutation (directive retirée d'une page de \code{ai-chatbot}).
\end{decision}

La direction de l'erreur est la bonne. Une arête manquante \emph{retire} des modules du graphe
serveur, donc des findings : la règle sous-déclare (\issue{29}), elle ne sur-déclare pas. Et comme
aucun module n'est sauté, cette sous-déclaration ne cache rien aux autres règles.

\section{Pièges et limites connues}
\label{sec:regles-rendu-rsc-limites}

\begin{piege}
\textbf{Le compte des rendus n'est pas toujours une borne inférieure.} \code{subtree_renders}
compte \emph{un} rendu pour un élément non résolu \emph{sous} un candidat
(\code{None => renders += 1}, \cref{lst:rendu-subtree}). Dans \file{/tmp/ch21/unres.tsx},
\code{function Tree() \{ return <div><Lib /></div>; \}} avec \code{Lib} importé de
\code{some-lib} donne \og re-renders \code{<Tree>} (2 component renders)\fg{}. Si \code{Lib} est
un \code{memo}, React ne le re-rend pas, et le vrai compte est 1. L'affirmation d'absence n'est
pas en cause (\code{Lib} ne reçoit rien) ; seul le seuil \code{minWastedRenders} peut être franchi
à tort. Au premier niveau, en revanche, le \code{<Lib />} du même fichier n'est jamais candidat.
\end{piege}

\begin{piege}
\textbf{Un \code{.map} ne fait pas toujours une liste.} Seul un site de \emph{composant} construit
dans un callback est \code{in_list}. Un \code{.map} qui ne produit que des éléments hôtes
(\code{<span>}, \code{<tr>}) ne rend pas le sous-arbre \og liste\fg{}, et son coût reste d'un
rendu (\cref{lst:rendu-hover}).
\end{piege}

\paragraph{Écarts de texte observés.} Trois défauts de gabarit, sans effet sur les décisions :
\og (1 component renders)\fg{} au pluriel dans le message principal
(\src{src/rules/impls/wasted_subtree_render.rs}{212}{216}), alors que la note \code{Step::Rerender}
dit \og 1 component render\fg{} ; \og each \code{setinterval} event\fg{} pour une minuterie ;
\og Server Component. this file…\fg{} en minuscule après un point
(\src{src/rules/impls/server_component_hook.rs}{179}{187}). Le commentaire de
\code{REACT_CLIENT_HOOKS} (\og The engine models the first five as dedicated \code{HookEntry}
kinds\fg{}) est ambigu : les cinq premiers de la liste (\code{useContext}\dots) n'ont pas de genre
dédié ; il faut lire \og les cinq hooks modélisés, traités par les bras dédiés de
\code{client_only_hook}\fg{}.

\paragraph{Limites connues de \regle{server-component-hook}.}
\begin{itemize}
  \item \textbf{Imports non résolus} (\issue{29}, \code{precision-fn}) : un sous-graphe atteint par
        un alias non résolu n'est pas classé.
  \item \textbf{Run rétréci} : \code{reactant check tests/fixtures/next_project/src/app} n'abaisse
        pas \file{counter.tsx}, seul porteur de \code{"use client"}. La porte 1 ferme la règle :
        aucun finding, mais le blind spot \og 3 imported file(s) resolved outside the analysed set
        and were never read\fg{} retient le \og clean bill\fg{} (\cref{chap:projet-cli-driver}).
  \item \textbf{Listes fermées} : un hook client d'un paquet non listé (\code{next-auth/react},
        \cref{ex:rsc-preuves}) n'est pas une preuve.
  \item \textbf{Hook en position \code{return}} : un hook personnalisé réduit à
        \code{return useMemo(…)} est vu, mais le finding sort sans numéro de ligne (\issue{140}).
  \item \textbf{Coût} : \code{server_modules} est recalculé à chaque appel, dans \code{check} et
        dans \code{safe_check}, pour chaque composant : $O(C \cdot (V+E))$ par exécution, là où
        le \code{RenderIndex} est mis en cache dans \code{ProgramCache}. Observation de coût, pas
        de justesse.
\end{itemize}

\paragraph{Limites connues de \regle{wasted-subtree-render}.} \code{memo} et \code{forwardRef}
opaques (\issue{64}) ; appel opaque supposé sans lecture ni écriture de module (\issue{51},
\issue{52}) ; une cascade à travers un contexte importé d'un paquet n'est pas suivie
(\file{docs/limitations.md}).

% ---------------------------------------------------------------------------
\begin{resume}
\begin{itemize}
  \item Les deux règles du chapitre lisent des faits calculés \emph{à côté} du domaine des
        valeurs : la dépendance de rendu (\adr{41}) pour l'une, les faits de module et une
        convention de noms (\adr{26}) pour l'autre. Toutes deux plafonnent à \Warning{}, l'une
        faute de défaut (un coût), l'autre faute de preuve dans le domaine.
  \item \regle{wasted-subtree-render} groupe les écritures par déclencheur (handler, ou
        enregistrements d'un effet), classe la fréquence (classement, jamais preuve), puis demande
        à \code{RenderIndex::wasted_siblings} quels éléments résolus aucune écriture du lot
        n'atteint, et compte leurs rendus (\code{subtree_renders}, liste comptée une fois).
  \item Son affirmation est une absence ; d'où l'\cref{inv:tout-inconnu-usage} : non résolu,
        récursion, profondeur, top, contexte innommable comptent comme usage. Le prix est une
        sous-déclaration assumée (\issue{64}).
  \item Deux options typées (\code{minWastedRenders}, \code{continuousOnly}) ajustent ce qui est
        montré, jamais ce qui est affirmé.
  \item \regle{server-component-hook} : porte \code{any_declares("use client")}, graphe
        \code{server_modules} (entrées App Router, atteignabilité arrêtée à la frontière), preuves
        de hook client en listes fermées, un finding par composant. Les modules serveur restent
        analysés par toutes les règles.
  \item Types : \code{RenderDeps}, \code{ElementSite}, \code{Relevance}, \code{Landing},
        \code{Wasted}, \code{TriggerKey}, \code{OptionSpec}, \code{ModuleFacts},
        \code{ModuleTable}. Fichiers :\\
        \file{src/rules/impls/wasted_subtree_render.rs},\\
        \file{src/rules/helpers/render_tree.rs}, \file{src/engine/render_deps.rs},\\
        \file{src/rules/impls/server_component_hook.rs}, \file{src/project/nextjs.rs}.
  \item Le chapitre suivant (\cref{chap:regles-declaratives}) quitte les règles natives pour les
        règles déclaratives des packs.
\end{itemize}
\end{resume}

\section{Exercices}
\label{sec:regles-rendu-rsc-exercices}

\begin{exercice}{Construire ou transmettre}{rendu-construire}
Dans \cref{lst:rendu-recherche}, proposez deux réécritures de \code{SearchPage} qui font taire
\regle{wasted-subtree-render}, l'une par extraction d'état, l'autre par \code{children}. Pour
chacune, dites quel résumé \code{RenderDeps} change et pourquoi \code{wasted_siblings} n'a plus
de candidat.
\end{exercice}
\begin{solution}
(1) Un composant \code{SearchBox} possède \code{query}, rend l'\code{<input>} et
\code{<Results>} ; \code{SearchPage} rend \code{<SearchBox />} et \code{<Sidebar />}. Le slot
appartient désormais à \code{SearchBox}, dont le seul site (\code{<Results>}) est touché.
\code{SearchPage} n'a plus d'état : aucun déclencheur. (2) \code{SearchPage(\{ sidebar \})} rend
\code{\{sidebar\}} et reçoit \code{<Sidebar />} de son parent : l'élément devient un site du
parent, et le résumé de \code{SearchPage} n'a plus de site \code{Sidebar} (comme
\file{children.tsx}).
\end{solution}

\begin{exercice}{Deux setters, un lot}{rendu-lot}
Un \code{onChange} appelle \code{setA(e.target.value)} puis \code{setB(true)}. Le composant rend
\code{<X a=\{a\} />} et \code{<Y b=\{b\} />}, plus \code{<Z />}. Combien de findings, et quels
éléments sont nommés ? Même question si \code{setB} est appelé depuis un \code{onBlur}.
\end{exercice}
\begin{solution}
Un seul déclencheur (même handler) : question \code{\{Slot(a), Slot(b)\}}, \code{X} et \code{Y}
touchés, seul \code{<Z>} est candidat ; le finding tire si \code{Z} coûte au moins
\code{minWastedRenders} rendus ou contient une liste. Avec \code{onBlur}, il y a deux
déclencheurs ; celui de \code{onBlur} est discret (\code{blur} n'est pas un événement de frappe)
et \code{b} est un booléen : muet par défaut. Celui du \code{onChange} a pour question
\code{\{Slot(a)\}} : \code{<Y>} et \code{<Z>} sont candidats. Rejoué (\file{/tmp/ch21/lot.tsx},
\code{Z} construisant deux \code{<Row>}) : \og each \code{change} event writes states \code{a}
and \code{b} and re-renders \code{<Z>} (at least 2 component renders, including a list)\fg{}
pour la première version, \og each \code{change} event writes state \code{a} and re-renders
\code{<Y>} and \code{<Z>} (at least 3 component renders, including a list)\fg{} pour la seconde.
\end{solution}

\begin{exercice}{La frontière}{rsc-frontiere}
Dans \cref{fig:rsc-graphe}, on ajoute \code{"use client"} en tête de \file{price.tsx}. Quels
modules restent serveur ? Que devient le finding sur \code{Price} ? Et si l'on retire la directive
de \file{search.tsx} au lieu d'en ajouter une ?
\end{exercice}
\begin{solution}
\file{price.tsx} devient une frontière : il n'est plus visité, et \code{server_modules} vaut
\{\file{app/page.tsx}, \file{app/shop/page.tsx}, \file{lib/format.ts}\}. Plus de finding sur
\code{Price} (rejoué : \og ✓ 6 file(s) no issues found.\fg{}). Sans directive dans \file{search.tsx}, le programme n'en contient plus aucune : la
porte 1 ferme la règle, aucun finding (c'est \file{/tmp/ch21/nx2/}).
\end{solution}

\begin{exercice}{Pourquoi pas une Error ?}{rsc-error}
Un contributeur propose une primitive must \code{proven_server_module(file)} qui rendrait un
\code{Certified} quand le fichier est une entrée \code{page} sans directive et sans aucun import.
Le cas semble certain. Qu'en dit \adr{26} ? Quel fait reste hors du domaine ?
\end{exercice}
\begin{solution}
Même sans import, le fait \og ce fichier est une entrée serveur\fg{} reste une convention de
chemin (\code{server_entry_kind}) : le routeur peut être configuré autrement, et le dossier
\file{app/} n'être pas celui de Next.js. \adr{26} §4 refuse précisément cette primitive, qui
\og habillerait une convention de chemin en preuve de domaine\fg{}.
\end{solution}
